\documentclass[11pt]{article}
\usepackage[margin=0.85in]{geometry}
\usepackage{amsmath,amssymb,booktabs,array,longtable,hyperref,xcolor,listings,microtype}
\usepackage[T1]{fontenc}
\usepackage{lmodern}
\hypersetup{colorlinks=true,linkcolor=blue!55!black,urlcolor=blue!55!black}
\definecolor{codebg}{RGB}{246,248,250}
\lstset{basicstyle=\ttfamily\small,backgroundcolor=\color{codebg},frame=single,
  breaklines=true,columns=fullflexible,keepspaces=true}
\newcommand{\code}[1]{\texttt{#1}}
\newcommand{\pathcode}[1]{\texttt{\detokenize{#1}}}
\title{GATED GLUE: Implementation Changes and Mathematical Conformance}
\author{Technical implementation report}
\date{Branch comparison: \code{master} \texttt{5a30146c} to
      \code{GATED\_GLUE} \texttt{26946eb}\\Compiled \today}

\begin{document}
\maketitle

\begin{abstract}
This report explains the static edge-gated extension to GLUE, derives its
forward propagation and regularized objective, and maps each mathematical
quantity to the exact implementation.  It also covers graph-topology identity,
pretrain/fine-tune transfer, serialization, diagnostics, tests, and the matched
10x benchmark.  The central claim is precise: GATED GLUE retains the original
signed, normalized GCN propagation, but multiplies every directed non-self-loop
edge by an independent sigmoid gate without renormalizing after gating.
The primary benchmark uses gate initialization $0.95$ and retention weight
$\lambda_{\mathrm{keep}}=0$, isolating the effect of trainable static gating.
\end{abstract}

\tableofcontents
\newpage

\section{Scope and branch-level summary}

The comparison contains six GATED-specific commits and changes 18 files
(1,210 insertions and 17 deletions).  Existing GCN behavior remains the default;
users must explicitly select \code{graph\_encoder="gated"}.  The changes divide
into four layers:
\begin{enumerate}
  \item a low-level signed edge-gated aggregation operator;
  \item a topology-aware variational graph encoder with trainable gate logits;
  \item model/trainer lifecycle support and optional retention regularization;
  \item evaluation, diagnostics, tests, and matched benchmark infrastructure.
\end{enumerate}

\begin{longtable}{p{0.31\linewidth}p{0.61\linewidth}}
\toprule
File & Responsibility \\
\midrule
\endhead
\pathcode{scglue/models/nn.py} & \code{EdgeGatedGraphConv}: gated signed aggregation. \\
\pathcode{scglue/models/sc.py} & \code{GatedGraphEncoder}: logits, topology identity, activation, keep loss, diagnostics. \\
\pathcode{scglue/models/glue.py} & Generic trainer integration of the keep penalty. \\
\pathcode{scglue/models/scglue.py} & Public model selection, graph configuration, transfer, validation, encoding. \\
\pathcode{scglue/models/__init__.py} & Graph configuration in the high-level two-stage fitting helper. \\
\pathcode{evaluation/workflow/scripts/run_GLUE.py} & CLI selection, two-stage GATED fitting, saved gate diagnostics. \\
\pathcode{tests/models/test_gated_graph.py} & Unit and lifecycle conformance tests. \\
\bottomrule
\end{longtable}

\section{Reference GCN mathematics}

Let $G=(V,E)$ be the directed graph representation used by GLUE.  Each edge
$e=(i\!\to\!j)$ has a sign $s_e$ and a precomputed normalization coefficient
$n_e$.  Let $h_i\in\mathbb{R}^d$ be the source vertex representation.  The
original GCN aggregation at target $j$ is
\begin{equation}
  m_j^{\mathrm{GCN}} = \sum_{e=(i\to j)\in E} s_e n_e h_i.
  \label{eq:gcn}
\end{equation}
The normalization coefficient is supplied to the layer as \code{enorm}; the
edge sign is \code{esgn}.  The GATED branch does not redefine either quantity.
This is important for matched comparisons: graph construction, signed edges,
and the original normalization all remain unchanged.

\section{Static edge-gated propagation}

\subsection{Gate parameterization}

For every directed non-self-loop edge $e$, the model introduces an unconstrained
trainable logit $a_e\in\mathbb{R}$ and activates it with a sigmoid:
\begin{equation}
  g_e = \sigma(a_e)=\frac{1}{1+\exp(-a_e)}, \qquad 0<g_e<1.
  \label{eq:sigmoid}
\end{equation}
Self-loops deliberately have no trainable logit and use
\begin{equation}
  g_{(i\to i)}=1.
  \label{eq:selfloop}
\end{equation}
For requested initial probability $p_0$ (\code{gate\_init}), each non-self-loop
logit is initialized by the inverse sigmoid:
\begin{equation}
  a_e^{(0)}=\operatorname{logit}(p_0)
  =\log\frac{p_0}{1-p_0}.
  \label{eq:init}
\end{equation}
At the benchmark value $p_0=0.95$, $a_e^{(0)}\approx2.944439$ and every
non-self-loop edge initially retains 95\% of its original normalized message.

\subsection{Forward equation}

GATED propagation is
\begin{equation}
  m_j^{\mathrm{GATED}}
  =\sum_{e=(i\to j)\in E}s_e n_e g_e h_i.
  \label{eq:gated}
\end{equation}
The implementation in \code{EdgeGatedGraphConv.forward} computes
\begin{lstlisting}[language=Python]
message = input[sidx] * (esgn * enorm * gate).unsqueeze(1)
result.scatter_add_(0, tidx.unsqueeze(1).expand_as(message), message)
\end{lstlisting}
Thus the exact code correspondence is
\begin{center}
\begin{tabular}{ll}
\toprule
Mathematics & Tensor/code quantity \\
\midrule
$h_i$ & \code{input[sidx]} \\
$s_e$ & \code{esgn} \\
$n_e$ & \code{enorm} \\
$g_e$ & \code{gate} from \code{activated\_gates} \\
$\sum_{i\to j}$ & target-indexed \code{scatter\_add\_} \\
\bottomrule
\end{tabular}
\end{center}

\subsection{No post-gate renormalization}

Equation~\eqref{eq:gated} uses the original $n_e$ unchanged.  It does
\emph{not} replace it by a gate-dependent normalization such as
\begin{equation*}
  \widetilde n_e = \frac{n_e g_e}{\sum_{e':\,\operatorname{target}(e')=j}n_{e'}g_{e'}}.
\end{equation*}
Consequently, reducing one gate attenuates precisely that message and cannot
inflate neighboring messages through a changed denominator.  This is the
meaning of ``static edge gating'' here: gates are global learned parameters
attached to graph edges, not sample-dependent attention coefficients, and they
directly modulate the already-normalized graph operator.

\subsection{Variational output}

After gated aggregation, the encoder retains the original variational form:
\begin{align}
  \mu_j &= W_\mu m_j+b_\mu,\\
  \sigma_j &= \operatorname{softplus}(W_\sigma m_j+b_\sigma)+\varepsilon,\\
  q(v_j\mid G)&=\mathcal{N}(\mu_j,\operatorname{diag}(\sigma_j^2)).
\end{align}
In code these are \code{loc}, \code{std\_lin}, and the returned
\code{torch.distributions.Normal}.  Therefore gating changes graph message
strengths but not the latent distribution family or graph decoder.

\section{Training objective and retention regularization}

Let $\mathcal L_{\mathrm{data}}$ denote the weighted cell-domain ELBO term,
$\mathcal L_{\mathrm{graph}}$ the graph ELBO term, and
$\mathcal L_{\mathrm{dsc}}$ the adversarial discriminator loss.  Omitting
supervised/paired extensions for clarity, the existing generator objective is
organized as
\begin{align}
  \mathcal L_{\mathrm{VAE}}^{\mathrm{base}}
    &=\lambda_{\mathrm{data}}\mathcal L_{\mathrm{data}}
      +\lambda_{\mathrm{graph}}D\mathcal L_{\mathrm{graph}},\\
  \mathcal L_{\mathrm{gen}}^{\mathrm{base}}
    &=\mathcal L_{\mathrm{VAE}}^{\mathrm{base}}
      -\lambda_{\mathrm{align}}\mathcal L_{\mathrm{dsc}},
\end{align}
where $D$ is the number of data domains.

The branch defines the raw gate-retention penalty over trainable non-self-loop
edges:
\begin{equation}
  \mathcal L_{\mathrm{keep}}=\sum_{e\in E_{\mathrm{nonself}}}(1-g_e)^2.
  \label{eq:keep}
\end{equation}
The modified objective is
\begin{align}
  \mathcal L_{\mathrm{VAE}}^{\mathrm{GATED}}
    &=\mathcal L_{\mathrm{VAE}}^{\mathrm{base}}
      +\lambda_{\mathrm{keep}}\mathcal L_{\mathrm{keep}},\\
  \mathcal L_{\mathrm{gen}}^{\mathrm{GATED}}
    &=\mathcal L_{\mathrm{VAE}}^{\mathrm{GATED}}
      -\lambda_{\mathrm{align}}\mathcal L_{\mathrm{dsc}}.
  \label{eq:objective}
\end{align}
This term is added consistently to generic GLUE, SCGLUE, independent SCGLUE,
and paired SCGLUE trainers.  The code computes
\code{((1 - sigmoid(gate\_logits)) ** 2).sum()} and then adds
\code{self.lam\_keep * gate\_keep\_loss} to \code{vae\_loss}.

The sign of the regularizer is correct for a retention prior: minimizing it
pushes $g_e$ toward one.  It does not promote sparsity.  The derivative with
respect to a logit is
\begin{equation}
  \frac{\partial\mathcal L_{\mathrm{keep}}}{\partial a_e}
  =-2(1-g_e)g_e(1-g_e),
\end{equation}
which is non-positive for $0<g_e<1$ and therefore increases the gate under
gradient descent.

For the primary matched benchmark,
\begin{equation}
  \boxed{\lambda_{\mathrm{keep}}=0},
\end{equation}
so Equation~\eqref{eq:objective} reduces to the original GLUE objective while
the graph operator remains Equation~\eqref{eq:gated}.  This isolates the
architectural effect of learned edge attenuation from retention regularization.

\section{Topology identity and invariants}

Independent gates are meaningful only if logit $a_e$ remains attached to the
same directed edge.  The implementation therefore persists the configured
edge index tensor and enforces the following invariants:
\begin{enumerate}
  \item edge indices have shape $2\times |E|$ and valid vertex bounds;
  \item duplicate directed edges are rejected;
  \item reconfiguration is allowed only for the same directed topology;
  \item reordered but equivalent edge lists are aligned by integer edge IDs;
  \item changed topology is rejected during compilation, adoption, fitting,
        loss evaluation, and graph encoding.
\end{enumerate}
For $|V|=N$, directed edge $(i,j)$ is represented by the collision-free ID
$iN+j$.  Sorting these IDs yields the permutation between equivalent edge
orders.  This makes the gate vector order-invariant without weakening topology
validation.

Graph configuration occurs before compilation because \code{gate\_logits} is
created when the topology is known.  Creating it after compilation would omit
the parameter from the optimizer.  The branch explicitly rejects compilation
of an unconfigured gated model and tests that gate logits occur in the VAE
optimizer parameter groups.

\section{Pretraining, fine-tuning, and serialization}

The existing GLUE workflow trains two models:
\begin{enumerate}
  \item pretraining with alignment burn-in disabled via an infinite burn-in;
  \item fine-tuning after estimating balancing weights from pretrained cell
        embeddings.
\end{enumerate}
Both GATED models are configured with the same graph before compilation.
During \code{adopt\_pretrained\_model}, architecture equality and topology
equivalence are checked, then learned gate logits are copied using edge
identity rather than raw array position.  This ensures that pretraining gates
survive into fine-tuning even if equivalent graphs enumerate edges differently.

The configured edge tensor, non-self-loop mask, and gate logits are included in
model state.  Save/load tests verify exact recovery.  Cross-architecture
adoption (GCN to GATED or GATED to GCN) is rejected rather than partially
copying an incompatible graph encoder.

\section{Backward compatibility}

The API defaults are
\begin{lstlisting}[language=Python]
graph_encoder = "gcn"
gate_init = 0.95
lam_keep = 0.0
\end{lstlisting}
When \code{graph\_encoder="gcn"}, the original \code{GraphEncoder} and
\code{GraphConv} are constructed.  The trainer's \code{gate\_keep\_loss}
returns a scalar zero for encoders without gates.  Therefore the new objective
term is exactly zero for GCN, and existing calls preserve their architecture
and numerical path except for graph-configuration calls that are no-ops for
GCN.

\section{Evaluation command and diagnostics}

The evaluation runner adds three validated options:
\begin{lstlisting}
--graph-encoder {gcn,gated}
--gate-init PROBABILITY_IN_(0,1)
--lam-keep FINITE_NONNEGATIVE_FLOAT
\end{lstlisting}
The same values are used in both training stages.  GATED runs save
\code{gate\_diagnostics.yaml} with
\begin{itemize}
  \item number of trainable non-self-loop gates;
  \item initial gate mean;
  \item final gate minimum, mean, and maximum;
  \item maximum absolute change in gate logits.
\end{itemize}
These are architectural diagnostics and are not mixed into biological or
integration scores.

\section{Mathematical conformance tests}

The new tests provide executable checks of the equations:
\begin{longtable}{p{0.34\linewidth}p{0.57\linewidth}}
\toprule
Tested property & Mathematical guarantee \\
\midrule
\endhead
Signed weighted example & Directly verifies Equation~\eqref{eq:gated}. \\
Fixed denominator after gate change & Confirms no post-gate renormalization. \\
Initialization at $0.95$ & Verifies Equation~\eqref{eq:init}. \\
Self-loop fixed at one & Verifies Equation~\eqref{eq:selfloop}. \\
Finite nonzero logit gradients & Gates participate in differentiable training. \\
Keep loss numeric example & Verifies Equation~\eqref{eq:keep}. \\
Optimizer membership & Topology is configured before optimizer construction. \\
Reordered edge adoption & Gate identity follows directed edges, not list position. \\
Changed/duplicate topology rejection & Prevents silent gate-to-edge corruption. \\
Serialization round trip & Preserves topology and gate values. \\
GCN unchanged/default & Preserves backward-compatible model selection. \\
CLI validation & Rejects invalid probabilities and negative keep weights. \\
\bottomrule
\end{longtable}

\section{Matched 10x benchmark design}

The branch adds nine GATED subsample runs for sizes 250, 2,000, and 8,000 with
subsample seeds 0, 1, and 2.  They reuse the exact prepared GCN/GAT cell subsets,
validated by ordered cell-ID hashes.  Only the graph encoder changes.  The
matched configuration is
\begin{center}
\begin{tabular}{lr}
\toprule
Parameter & Value \\
\midrule
Learning rate & 0.002 \\
$\lambda_{\mathrm{graph}}$ & 0.02 \\
$\lambda_{\mathrm{align}}$ & 0.05 \\
Latent / alternative dimension & 50 / 100 \\
Hidden depth / dimension & 2 / 256 \\
Dropout & 0.2 \\
Negative samples per edge & 10 \\
Data batch size & 128 \\
Graph encoder & gated \\
Gate initialization & 0.95 \\
$\lambda_{\mathrm{keep}}$ & 0.0 \\
Model seed & 0 \\
\bottomrule
\end{tabular}
\end{center}
Preprocessing, guidance graph construction, two-stage fitting, automatic graph
batching, balancing weights, convergence requirements, and early stopping are
the same as the matched benchmark.

Feature consistency requires an architecture-specific full-data feature
embedding.  The branch therefore adds one full-data GATED seed-0 reference.
For every subsample, the metric compares its GATED \code{feature\_latent.csv}
against the full GATED \code{feature\_latent.csv}; it explicitly rejects a
symlinked substitute.  GCN and GAT retain their own full architecture-specific
references.

\section{Interpretation and limitations}

The learned $g_e$ should be interpreted as a global edge-retention multiplier,
not a probability that an edge is sampled, a causal importance score, or a
cell-specific attention coefficient.  Gates can attenuate but cannot reverse
an edge sign or amplify it above its original normalized strength.  Because
self-loops remain one, each vertex always retains its original self-message.

With $\lambda_{\mathrm{keep}}=0$, gates are controlled only by the original
training objective.  This is appropriate for isolating the gating architecture,
but it does not impose sparsity or guarantee gates remain near one.  Future
retention experiments may vary $\lambda_{\mathrm{keep}}$, but those constitute
a different comparison and should not replace the primary matched benchmark.

\section{Conclusion}

The implementation conforms to the specified mathematics at each layer:
sigmoid logits provide bounded independent non-self-loop gates;
Equation~\eqref{eq:gated} is implemented as direct signed normalized message
multiplication with no renormalization; Equation~\eqref{eq:keep} is included
with the correct sign and weight; topology-aware lifecycle code preserves the
edge-to-gate mapping across training stages and serialization; and the benchmark
changes only the intended architecture-specific quantities.  The default GCN
path remains available and is the default model selection.

\appendix
\section{Relevant implementation code}

This appendix reproduces the relevant source directly from the compared
\code{GATED\_GLUE} branch.  The displayed line numbers are repository source
line numbers, allowing every equation in the main report to be audited against
the implementation.

\subsection{Signed gated aggregation}

This is the complete low-level operator implementing
$s_e n_e g_e h_i$ followed by target-indexed summation.  Notice that
\code{enorm} is multiplied by \code{gate} and is never recomputed.

\lstinputlisting[language=Python,numbers=left,firstnumber=57,firstline=57,lastline=87]{scglue/models/nn.py}

\subsection{Gate construction, topology identity, activation, and encoder output}

The following listing contains the trainable-logit construction, inverse
sigmoid initialization, fixed self-loops, edge-ID permutation, retention loss,
gate accessor, and variational forward pass.

\lstinputlisting[language=Python,numbers=left,firstnumber=52,firstline=52,lastline=183]{scglue/models/sc.py}

\subsection{Retention penalty in the GLUE objective}

The generic trainer returns zero for a non-gated encoder and otherwise obtains
the raw penalty from the gated encoder.  Lines 422--426 show its exact placement
in the VAE and generator objectives.

\lstinputlisting[language=Python,numbers=left,firstnumber=346,firstline=346,lastline=350]{scglue/models/glue.py}
\lstinputlisting[language=Python,numbers=left,firstnumber=390,firstline=390,lastline=431]{scglue/models/glue.py}

\subsection{Public model selection and graph configuration}

The constructor preserves GCN as the default and selects GATED explicitly.
The graph configuration method converts the full graph into the exact directed
edge representation used during fitting before creating the optimizer.

\lstinputlisting[language=Python,numbers=left,firstnumber=1278,firstline=1278,lastline=1302]{scglue/models/scglue.py}
\lstinputlisting[language=Python,numbers=left,firstnumber=1357,firstline=1357,lastline=1385]{scglue/models/scglue.py}

\subsection{Compilation guard and objective wiring}

This code prevents gated compilation before graph configuration and passes the
validated \code{lam\_keep} value into the trainer.

\lstinputlisting[language=Python,numbers=left,firstnumber=1463,firstline=1463,lastline=1511]{scglue/models/scglue.py}

\subsection{Two-stage evaluation workflow}

The evaluation workflow uses identical gated settings in pretraining and
fine-tuning, configures the graph before each compilation, and transfers the
pretrained model before the second optimizer is constructed.

\lstinputlisting[language=Python,numbers=left,firstnumber=205,firstline=205,lastline=272]{evaluation/workflow/scripts/run_GLUE.py}

\subsection{Saved gate diagnostics}

The final model exposes its activated gates and the workflow saves architectural
diagnostics separately from integration metrics.

\lstinputlisting[language=Python,numbers=left,firstnumber=274,firstline=274,lastline=317]{evaluation/workflow/scripts/run_GLUE.py}

\subsection{Executable mathematical checks}

These tests numerically verify signed gated propagation, the fixed normalization
denominator, inverse-sigmoid initialization, fixed self-loops, the keep loss,
and differentiable gate updates.

\lstinputlisting[language=Python,numbers=left,firstnumber=18,firstline=18,lastline=78]{tests/models/test_gated_graph.py}

The lifecycle test below verifies optimizer membership, pretrain/fine-tune
adoption, edge-order alignment, serialization, and rejection of incompatible
topologies or architectures.

\lstinputlisting[language=Python,numbers=left,firstnumber=136,firstline=136,lastline=205]{tests/models/test_gated_graph.py}

\end{document}
